\section{Methods}
\label{sec:methods}

\subsection{The scoring system}
\textsc{AIPR} grades a manuscript in a single frozen pipeline version (v6), applying the
now-standard practice of scoring by prompting an LLM against a rubric
\citep{liu2023geval,zheng2023judging} to a new object: the manuscript. The pipeline
uploads the submitted PDF and runs two passes. A \emph{reviewer} pass reads the paper
and emits the five dimension scores with justifications and verbatim-anchored
highlights. An \emph{audit} pass re-reads the paper with a live literature-search tool,
checks each highlight against the manuscript, and proposes missing citations against a
real bibliographic index, so a recommended citation is always a record the tool
returned, not a generated string. The overall score is a fixed weighted mean of the five
dimensions,
\begin{equation}
  \text{overall} = \frac{4\,s_{\text{nov}} + 2\,s_{\text{rig}}
  + 4\,s_{\text{app}} + 1\,s_{\text{cla}} + 0.5\,s_{\text{cit}}}{11.5},
  \label{eq:overall}
\end{equation}
with novelty and applicability weighted most and citation least; the model does not emit
the overall directly. The pipeline's product is this structured, evidence-grounded
review; the overall score is a by-product of producing it, not the system's purpose, and
the baseline of \S\ref{sec:res-value} returns only the number. The system is prompt-only:
it is not fine-tuned on reviews, decisions, or ratings, and no model here is trained on
the evaluation data.

The contrast with the baseline is concrete. Where the baseline prompt
(\S\ref{sec:res-value}) is one paragraph of about fifty words, the reviewer pass is
driven by a ${\approx}2{,}200$-word rubric instruction paired with a $1{,}200$-word
output schema, and the audit pass adds ${\approx}1{,}200$ words plus a tool loop of up to
eight OpenAlex search rounds. The two passes are role-separated, opening respectively
``You are an expert peer reviewer\,\ldots\ Produce a structured review'' and ``You are a
senior scientific editor\,\ldots\ audit the reviewer's highlights against the paper, and
find real missing citations.'' The full rubric, anchors, and schema are proprietary
(\S\ref{sec:discussion}); the point of \S\ref{sec:res-value} is that this instruction
footprint, almost two orders of magnitude larger than the baseline, buys reliability and
a grounded review rather than a higher discrimination score.

\subsection{Configurations and the baseline}
Two configurations run the identical two-pass pipeline and differ only in the model:
\textbf{AIPR (GPT-5.4-mini)} on the cheaper model and \textbf{AIPR (GPT-5.4)} on the
frontier model, the production configuration.\footnote{Exported as \texttt{full\_mini}
and \texttt{full\_full} in the released \texttt{gradings.csv}; we use the display names
throughout. The baseline \textbf{Direct (GPT-5.4)} is exported as \texttt{naive}.} The
cheap configuration is effectively unconstrained in volume; the frontier configuration is
budgeted. We therefore establish the score--outcome relationship on the large cohort $M$
and confirm it on cohort $H$, after verifying the two scores agree well enough for the
former to proxy the latter (H5). The pre-registered baseline, \textbf{Direct (GPT-5.4)},
runs the same frontier model on the same PDF with a single one-paragraph prompt and no
rubric or audit. We treat the system as a fixed black box: the dimension weights were set
by empirical internal validation on a separate, held-out manuscript set during product
development (no ICLR submissions from any year, none of the papers, decisions, or ratings
used here) and frozen before this study, not fit to any outcome in this cohort. The
lightly weighted citation dimension was uninformative in this run for a technical reason
in its audit, so we exclude it from interpretation; the headline is unchanged when it is
dropped, and a grounded, search-backed citation signal is an obvious next step
(\S\ref{sec:discussion}). We evaluate the rule \emph{as deployed} and report weighting
sensitivity (equal-weight, leave-one-out) in Appendix~\ref{app:robust}.

\subsection{Hypotheses, metrics, and pre-registration}
All hypotheses, the primary metric, and the low-score threshold were fixed before any
score was joined to any outcome (Appendix~\ref{app:prereg}). We report effect sizes with
bootstrap confidence intervals throughout, because ``strong signal'' is an effect-size
claim rather than a $p$-value claim. H1--H4 assemble criterion-related validity evidence
for a specific interpretation (weak relative to the venue bar) under a specific use
(triage), following the \emph{Standards for Educational and Psychological Testing}
\citep{aera2014standards}; a separate comparison (V1) asks whether that validity is the
pipeline's or the model's.

\begin{description}\itemsep2pt
  \item[H1 (primary): low-end flagging.] Submissions in the lowest score quintile are
  rejected well above the base rate, with oral papers essentially absent. Reported as
  reject rate and \emph{lift} over the base rate, with Wilson intervals.
  \item[H2: separation.] The overall score discriminates reject from accept (AUROC,
  stratified bootstrap CI).
  \item[H3: monotone gradient.] Mean overall score increases across the three ordered
  tiers (Jonckheere--Terpstra trend test, Monte Carlo permutation $p$; Spearman $\rho$).
  \item[H4: continuous agreement.] The overall score correlates with the mean reviewer
  rating (Spearman $\rho$, bootstrap CI).
  \item[H5 (gate): mini-to-frontier bridge.] On cohort $H$, AIPR (GPT-5.4-mini) and
  AIPR (GPT-5.4) overall scores correlate at $\rho \ge 0.8$, licensing the large-$N$
  cohort-$M$ results as a proxy for the production score.
  \item[V1 (primary value): pipeline vs.\ baseline.] On cohort $H$, AIPR (GPT-5.4)
  out-discriminates Direct (GPT-5.4) on the same model and PDF. Pre-declared success:
  AUROC difference $>0$ with the CI excluding zero; reported alongside accept/reject
  agreement and run-to-run reliability. A null is reported plainly.
\end{description}
Confidence intervals are 4{,}000-resample BCa bootstraps (stratified for class-conditional
statistics such as AUROC; Appendix~\ref{app:power}); the trend test uses 10{,}000 Monte
Carlo permutations; multiple comparisons across the five dimensions use Benjamini--Hochberg;
and the analysis is deterministic under a fixed seed. Secondary analyses (per-dimension
discrimination, audit-pass contribution, run-to-run variance, prevalence re-weighting,
weighting robustness, leakage controls) are reported in \S\ref{sec:results} and
Appendix~\ref{app:robust}.

\paragraph{Deviations from pre-registration.}
Two estimator choices differ from the registered analysis plan. Both are
analysis-stage substitutions of a better-behaved interval for the registered
one, and neither alters a pre-declared success threshold. First, confidence
intervals are BCa bootstraps rather than the registered percentile bootstrap:
the estimator-validation simulation (Appendix~\ref{app:power}) shows the
percentile interval under-covers at this sample size while BCa attains nominal
coverage. The V1 paired AUROC-difference CI is the exception and remains a
percentile interval on the paired bootstrap distribution. Second, the H1
success rule (lift CI lower bound above 1) reads a whole-cohort bootstrap CI
for the bottom-band lift rather than the registered Wilson-based interval:
resampling the whole cohort propagates uncertainty in the base rate and in the
sample-relative band edges, which a Wilson interval on the band's reject rate
alone cannot capture. Wilson intervals are still reported descriptively
alongside the bands (Table~\ref{tab:bands}).
